\section{Where do Lie algebras come from?}

\begin{definition}
    A Lie algebra is a vector space $ L $ over a field $ F $, with a binary operation $ [a,b] \in L $, such that
    \begin{enumerate}
        \item $ [a,b] $ is bilinear.
        \item $ [a,b] = -[b,a] $.
        \item $ [a,[b,c]] + [c,[a,b]] + [b,[c,a]] = 0 $, this is called Jacobi identity.
    \end{enumerate}
\end{definition}

\subsection{They come from associative algebras}

Let $ A $ be an associative algebra. $ \forall a,b \in A $, define $ [a,b]=ab-ba $. Check that $ (A,[ \ ,\  ]) $ is a Lie algebra:

\begin{proof}
    Bilinearity and skew-symmetry are clear. For Jacobi identity:
    \begin{align*}
        [a,[b,c]] & + [b,[c,a]] + [c,[a,b]]                                                 \\
                  & = a(bc-cb) - (bc-cb)a + b(ca-ac) - (ca-ac)b + c(ab-ba) - (ab-ba)c       \\
                  & = abc - acb - bca + cba + bca - bac - cab + acb + cab - cba - abc + bac \\
                  & = 0.
    \end{align*}
\end{proof}

We denote this new algebra $ (A,[ \ ,\  ]) $ by $ A^{(-)} $.

Suppose that $ L \subseteq A $ and $ [L,L] \subseteq L $. Then $ L \le A^{(-)} $。

\begin{example}
    $ M_n(F)^{(-)}=\mathfrak{gl}(n) $. We call it the general linear Lie algebra.

    $ [L,L]=\{ \sum_{i} [a_i,b_i] \} \trianglelefteq L $. Clearly $ [\mathfrak{gl}(n),\mathfrak{gl}(n)] \neq \mathfrak{gl}(n) $ and for $ n \ge 2 $, $ [\mathfrak{gl}(n),\mathfrak{gl}(n)]\neq 0 $.
\end{example}

\begin{example}
    Focus on $ [\mathfrak{gl}(n),\mathfrak{gl}(n)] $, all matrices are of trace 0. This is called the special linear Lie algebra, denoted by $ \mathfrak{sl}(n) $.
\end{example}

\begin{example}
    If the base field changed, some properties of the Lie algebra may change. For example, if $\operatorname{char}F=2$, then $[a,b]=[b,a]$. The Lie bracket is no longer skew-symmetric but actually symmetric.

    Consider $ [\mathfrak{gl}(n),\mathfrak{gl}(n)]=\mathfrak{sl}(n) $. If $\operatorname{char}F=0$ or $\operatorname{char}F \nmid n$, then $\mathfrak{sl}(n)$ is a simple Lie algebra. If $\operatorname{char}F\mid n$, then $\mathfrak{sl}(n)$ is not simple.
\end{example}

\begin{definition}
    $ Z $ is the center of Lie algebra $ L $ if $ Z=\{ z \in L | [z,L]=0 \} $.
\end{definition}

\subsection{They come from involutions of central simple algebras}

We can also find Lie algebras from the involution of a central simple algebra. Let $ A $ be central simple algebra and $ *: A \to A $ be a linear transformation. $ * $ is called a an involution if $ (a^*)^* =a $ and $ (ab)^*=b^*a^* $, $ \forall a,b \in A $.

\begin{proposition}
    Let $A$ be a central simple algebra over a field $F$ equipped with an involution
    $* : A \to A$. Then the set
    \[
        L=\{x\in A \mid *(x)=-x\}
    \]
    forms a Lie algebra under the commutator bracket
    \[
        [x,y]=xy-yx .
    \]
\end{proposition}

\begin{proof}
    First observe that any associative algebra $A$ becomes a Lie algebra
    with the commutator bracket
    \[
        [x,y]=xy-yx .
    \]

    Let
    \[
        L=\{x\in A\mid *(x)=-x\}.
    \]
    We show that $L$ is closed under the Lie bracket.

    Take $x,y\in L$. Then
    \[
        *(x)=-x,\qquad *(y)=-y.
    \]

    Since $*$ is an involution, it satisfies
    \[
        *(xy)=*(y)*(x).
    \]

    Hence
    \[
        *([x,y])
        =*(xy-yx)
        =*(y)*(x)-*(x)*(y).
    \]

    Substituting $*(x)=-x$ and $*(y)=-y$, we obtain
    \[
        *([x,y])
        =(-y)(-x)-(-x)(-y)
        =yx-xy
        =-[x,y].
    \]

    Thus $[x,y]\in L$. Therefore $L$ is closed under the commutator bracket.

    Since the commutator bracket in any associative algebra satisfies
    bilinearity, antisymmetry, and the Jacobi identity,
    it follows that $L$ is a Lie algebra.
\end{proof}

This construction explains why many classical Lie algebras arise from
central simple algebras with involution.  In particular, the classical
Lie algebras of Cartan--Killing type can be obtained in this way.

\begin{example}[Orthogonal involution]
    Let
    \[
        A=M_n(F)
    \]
    and define the involution
    \[
        *(X)=X^{T}.
    \]
    Then
    \[
        L=\{X\in M_n(F)\mid X^T=-X\}.
    \]
    This is the Lie algebra of skew-symmetric matrices, which is the
    orthogonal Lie algebra
    \[
        \mathfrak{so}_n(F).
    \]
\end{example}

\begin{example}[Symplectic involution]
    Let
    \[
        A=M_{2n}(F)
    \]
    and let
    \[
        J=
        \begin{pmatrix}
            0    & I_n \\
            -I_n & 0
        \end{pmatrix}.
    \]
    Define
    \[
        *(X)=J^{-1}X^{T}J.
    \]
    Then
    \[
        L=\{X\in M_{2n}(F)\mid *(X)=-X\}
    \]
    is the symplectic Lie algebra
    \[
        \mathfrak{sp}_{2n}(F).
    \]
\end{example}

\begin{example}[Unitary involution]
    Let $K/F$ be a quadratic field extension with Galois involution
    $a\mapsto \overline{a}$.  Consider
    \[
        A=M_n(K)
    \]
    with involution
    \[
        *(X)=\overline{X}^{\,T}.
    \]
    Then
    \[
        L=\{X\in M_n(K)\mid *(X)=-X\}
    \]
    is the unitary Lie algebra
    \[
        \mathfrak{su}_n .
    \]
\end{example}

Therefore classical Lie algebras arise from pairs $(A,*)$ where
$A$ is a central simple algebra and $*$ is an involution.  This
provides an algebraic framework for the Cartan--Killing classification
of classical Lie algebras.

\begin{example}
    The Quaternion algebra is an associative algebra with the involution $ a \mapsto a^* $. It is 4 dimensional central simple algebra with standard basis $ 1,i,j,k $, where the relations are $ i^2 = j^2 = k^2  = -1, ij=k, jk=i, ki=j $ as the following graph:
    \[
        \begin{tikzcd}
            i \arrow[r] & j \arrow[d]  \\
            & k \arrow[lu]
        \end{tikzcd}
    \]
    We use $ \mathbb{H} $ to denote the Quaternion algebra when we choose $ \mathbb{R} $ as the base field. i.e. $ \mathbb{H}=\{ a_0 1 + a_1 i + a_2 j + a_3 k | a_0,a_1,a_2,a_3 \in \mathbb{R} \} $. $ \forall x \in \mathbb{H} $ where $ x=a_0 1 + a_1 i + a_2 j + a_3 k $, $ * $ is defined as $ x^* = a_0 1 - a_1 i - a_2 j - a_3 k $ and $ x^{-1}= \frac{1}{a_0^2 + a_1^2 + a_2^2 + a_3^2} (a_0 1 - a_1 i - a_2 j - a_3 k) $. Hence $ \mathbb{H} $ is a division algebra.

    If the base field is algebraically closed, $F =\overline{F}$. We can denote the Quaternion algebra by $ 2 \times 2 $ matrices
    \[
        \begin{pmatrix}
            \alpha & \beta \\
            \gamma & \xi   \\
        \end{pmatrix}
    \]
    Where $ \alpha, \beta, \gamma, \xi \in F $. In this case it is no longer a division algebra ($a_0^2 + a_1^2 + a_2^2 + a_3^2$ may equal to 0), but just $ M_2(F) $. And in this case the involution $ * $ is defined as:
    \[
        \begin{pmatrix}
            \alpha & \beta \\
            \gamma & \xi   \\
        \end{pmatrix}
        \mapsto
        \begin{pmatrix}
            \xi     & -\beta \\
            -\gamma & \alpha \\
        \end{pmatrix}
    \]
    Which is the Symplectic involution.
\end{example}

\begin{example}
    From what we discussed. For a $2n \times 2n$ matrix $ A =\begin{pmatrix} A_1 & A_2 \\ A_3 & A_4 \end{pmatrix}$ in which $ A_1, A_2, A_3, A_4 $ are $ n \times n $ matrices, the unitary involution is defined as:
    \[
        A^*=
        \begin{pmatrix}
            A_4^t  & -A_3^t \\
            -A_2^t & A_1^t  \\
        \end{pmatrix}
    \]
    And the unitary involution is defined as:
    \[
        A^*=\overline{A}^t
    \]
    When $n=1$ we have two different type, and we will get two different Lie algebras. But they are isomorphic.
\end{example}

\begin{proposition}
    Let $F$ be a field of characteristic $\neq 2$, and let
    \[
        J=\begin{pmatrix}0&1\\-1&0\end{pmatrix}\in M_{2}(F).
    \]
    Define an involution $*$ on $M_{2}(F)$ by
    \[
        X^{*}=J^{-1}X^{t}J .
    \]
    Then the Lie algebra of $*$-skew elements,
    \[
        K(M_{2}(F),*)=\{X\in M_{2}(F)\mid X^{*}=-X\},
    \]
    is exactly $\mathfrak{sl}_{2}(F)$.
\end{proposition}

\begin{proof}
    Let
    \[
        X=\begin{pmatrix}a&b\\ c&d\end{pmatrix}\in M_{2}(F).
    \]
    Since
    \[
        J^{-1}=-J=\begin{pmatrix}0&-1\\1&0\end{pmatrix},
    \]
    we compute
    \[
        X^{*}
        =J^{-1}X^{t}J
        =\begin{pmatrix}0&-1\\1&0\end{pmatrix}
        \begin{pmatrix}a&c\\ b&d\end{pmatrix}
        \begin{pmatrix}0&1\\-1&0\end{pmatrix}
        =\begin{pmatrix}d&-b\\-c&a\end{pmatrix}.
    \]
    Now the condition $X^{*}=-X$ becomes
    \[
        \begin{pmatrix}d&-b\\-c&a\end{pmatrix}
        =
        \begin{pmatrix}-a&-b\\-c&-d\end{pmatrix}.
    \]
    Hence $d=-a$, while $b$ and $c$ are arbitrary. Therefore
    \[
        K(M_{2}(F),*)
        =
        \left\{
        \begin{pmatrix}
            a & b  \\
            c & -a
        \end{pmatrix}
        \,\middle|\,
        a,b,c\in F
        \right\}.
    \]
    But this is precisely the set of all $2\times 2$ matrices of trace $0$, namely
    \[
        \mathfrak{sl}_{2}(F)
        =
        \left\{
        \begin{pmatrix}
            a & b  \\
            c & -a
        \end{pmatrix}
        \,\middle|\,
        a,b,c\in F
        \right\}.
    \]
    Thus
    \[
        K(M_{2}(F),*)=\mathfrak{sl}_{2}(F).
    \]
    Since $K(M_{2}(F),*)$ is closed under the commutator bracket, it follows that the Lie algebra obtained from the symplectic involution on $M_{2}(F)$ is $\mathfrak{sl}_{2}(F)$.
\end{proof}

\begin{example}
    Let $K/F$ be a quadratic field extension, and let $a \mapsto \overline{a}$ denote the nontrivial
    $F$-automorphism of $K$. Consider the algebra $M_2(K)$ and define an involution $*$ by
    \[
        X^* = \overline{X}^{\,t}.
    \]
    Then $*$ is a unitary involution on $M_2(K)$.

    The Lie algebra of $*$-skew elements is
    \[
        K(M_2(K),*)=\{X\in M_2(K)\mid X^*=-X\}.
    \]
    If
    \[
        X=\begin{pmatrix} a & b \\ c & d \end{pmatrix},
    \]
    then
    \[
        X^*=\overline{X}^{\,t}
        =
        \begin{pmatrix}
            \overline{a} & \overline{c} \\
            \overline{b} & \overline{d}
        \end{pmatrix}.
    \]
    Thus the condition $X^*=-X$ is equivalent to
    \[
        \overline{a}=-a,\qquad \overline{d}=-d,\qquad \overline{c}=-b.
    \]
    Hence every element of $K(M_2(K),*)$ has the form
    \[
        X=
        \begin{pmatrix}
            a             & b \\
            -\overline{b} & d
        \end{pmatrix},
        \qquad \text{with } \overline{a}=-a,\ \overline{d}=-d.
    \]
    This is the unitary Lie algebra associated with the standard hermitian form.

    Its derived Lie algebra consists of those elements with trace $0$, namely
    \[
        \mathfrak{su}_2(K/F)
        =
        \left\{
        \begin{pmatrix}
            a             & b  \\
            -\overline{b} & -a
        \end{pmatrix}
        \,\middle|\,
        a,b\in K,\ \overline{a}=-a
        \right\}.
    \]
    After extending scalars to a splitting field (for instance to an algebraic closure of $F$),
    this Lie algebra becomes isomorphic to $\mathfrak{sl}_2$. Therefore the Lie algebra arising
    from a unitary involution in the $2\times 2$ case is of type $A_1$.
\end{example}

\begin{remark}
    Why does this Lie algebra become $\mathfrak{sl}_2$? The reason is that unitary Lie algebras are
    forms of special linear Lie algebras. In the present $2\times 2$ case, the corresponding type is
    $A_1$. Over a field where the quadratic extension splits, the involution becomes equivalent to the
    ordinary transpose situation, and the resulting special unitary Lie algebra becomes the split simple
    Lie algebra of type $A_1$, namely
    \[
        \mathfrak{sl}_2.
    \]
    Thus, although over the ground field one naturally writes the Lie algebra as $\mathfrak{su}_2$,
    its split form is precisely $\mathfrak{sl}_2$.
\end{remark}

\begin{example}[The unitary involution on $M_{2}(\mathbb{C})$ and the Lie algebra $\mathfrak{su}_{2}$]
    Let $F=\mathbb{R}$ and $K=\mathbb{C}$, and let $\overline{\phantom{x}}$ denote complex conjugation.
    Consider the algebra $M_{2}(\mathbb{C})$ equipped with the involution
    \[
        X^{*}=\overline{X}^{\,t}.
    \]
    This is the standard unitary involution on $M_{2}(\mathbb{C})$ relative to the extension
    $\mathbb{C}/\mathbb{R}$.

    We first determine the Lie algebra of $*$-skew elements:
    \[
        K(M_{2}(\mathbb{C}),*)=\{X\in M_{2}(\mathbb{C})\mid X^{*}=-X\}.
    \]
    Let
    \[
        X=\begin{pmatrix}a&b\\ c&d\end{pmatrix}\in M_{2}(\mathbb{C}).
    \]
    Then
    \[
        X^{*}
        =
        \overline{X}^{\,t}
        =
        \begin{pmatrix}
            \overline{a} & \overline{c} \\
            \overline{b} & \overline{d}
        \end{pmatrix}.
    \]
    Thus the condition $X^{*}=-X$ is equivalent to
    \[
        \overline{a}=-a,\qquad \overline{d}=-d,\qquad \overline{c}=-b.
    \]
    Hence $a$ and $d$ are purely imaginary, and $c=-\overline{b}$. Therefore
    \[
        K(M_{2}(\mathbb{C}),*)
        =
        \left\{
        \begin{pmatrix}
            a             & b \\
            -\overline{b} & d
        \end{pmatrix}
        \,\middle|\,
        a,d\in i\mathbb{R},\ b\in\mathbb{C}
        \right\}.
    \]
    This is the unitary Lie algebra $\mathfrak{u}_{2}$.

    Its derived Lie algebra is the trace-zero part, namely
    \[
        \mathfrak{su}_{2}
        =
        \left\{
        \begin{pmatrix}
            a             & b  \\
            -\overline{b} & -a
        \end{pmatrix}
        \,\middle|\,
        a\in i\mathbb{R},\ b\in\mathbb{C}
        \right\}.
    \]
    Indeed, $\operatorname{tr}([X,Y])=0$ for all $X,Y\in M_{2}(\mathbb{C})$, so every commutator in
    $\mathfrak{u}_{2}$ has trace $0$. Conversely, in this $2\times2$ case, the trace-zero subspace is
    generated by commutators, and hence it is exactly the derived Lie algebra.

    To see that $\mathfrak{su}_{2}$ is $3$-dimensional over $\mathbb{R}$, write
    \[
        a=it,\qquad b=u+iv,
        \qquad t,u,v\in\mathbb{R}.
    \]
    Then every element of $\mathfrak{su}_{2}$ has the form
    \[
        \begin{pmatrix}
            it    & u+iv \\
            -u+iv & -it
        \end{pmatrix},
        \qquad t,u,v\in\mathbb{R}.
    \]
    Therefore
    \[
        \dim_{\mathbb{R}}\mathfrak{su}_{2}=3.
    \]
    A convenient real basis is
    \[
        H=
        \begin{pmatrix}
            i & 0  \\
            0 & -i
        \end{pmatrix},
        \qquad
        X=
        \begin{pmatrix}
            0  & 1 \\
            -1 & 0
        \end{pmatrix},
        \qquad
        Y=
        \begin{pmatrix}
            0 & i \\
            i & 0
        \end{pmatrix}.
    \]
    One checks directly that
    \[
        [H,X]=2Y,\qquad [H,Y]=-2X,\qquad [X,Y]=2H,
    \]
    so $\mathfrak{su}_{2}$ is a $3$-dimensional real Lie algebra.

    Now we explain why this Lie algebra becomes $\mathfrak{sl}_{2}(\mathbb{C})$ after extending scalars
    from $\mathbb{R}$ to $\mathbb{C}$. Consider the complexification
    \[
        \mathfrak{su}_{2}\otimes_{\mathbb{R}}\mathbb{C}.
    \]
    Since $H,X,Y$ form a real basis of $\mathfrak{su}_{2}$, they also form a complex basis of
    $\mathfrak{su}_{2}\otimes_{\mathbb{R}}\mathbb{C}$. Define
    \[
        h=-iH=
        \begin{pmatrix}
            1 & 0  \\
            0 & -1
        \end{pmatrix},
        \
        e=\frac{1}{2}(X-iY)=
        \begin{pmatrix}
            0 & 1 \\
            0 & 0
        \end{pmatrix},
        \
        f=\frac{1}{2}(-X-iY)=
        \begin{pmatrix}
            0  & 0 \\
            -1 & 0
        \end{pmatrix}.
    \]
    These elements lie in $\mathfrak{su}_{2}\otimes_{\mathbb{R}}\mathbb{C}$ and satisfy
    \[
        [h,e]=2e,\qquad [h,f]=-2f,\qquad [e,f]=h.
    \]
    Thus they form a standard $\mathfrak{sl}_{2}$-triple, and therefore
    \[
        \mathfrak{su}_{2}\otimes_{\mathbb{R}}\mathbb{C}\cong \mathfrak{sl}_{2}(\mathbb{C}).
    \]

    In other words, $\mathfrak{su}_{2}$ is a real form of $\mathfrak{sl}_{2}(\mathbb{C})$. Hence the Lie
    algebra arising from the unitary involution on $M_{2}(\mathbb{C})$ is of type $A_{1}$.
\end{example}

\begin{remark}
    This explains why, in the $2\times 2$ case, the Lie algebra arising from a symplectic involution is isomorphic to the one of type $A_{1}$. Equivalently, in this lowest-rank case one has the accidental isomorphism
    \[
        C_{1}\cong A_{1}.
    \]
\end{remark}

\subsection{They come from the derivation of any algebra}

Let $A$ be any algebra.

\begin{definition}
    A derivation of $A$ is a linear map $\mathrm{d}: A \to A$ such that $\mathrm{d}(ab)=\mathrm{d}(a)b+a\mathrm{d}(b)$ for all $a,b \in A$.
\end{definition}

\begin{example}
    Let $A=F[t_1,\ldots ,t_n]$ be the polynomial algebra. We have the formal derivative $ \partial $ where $\partial(t_i^n)=nt_i^{n-1}$ and $\partial(fg)=f\partial(g)+\partial(f)g$.

    Let $\mathrm{d}=\sum_{i=1}^n f_i(t_1,\ldots ,t_n) \frac{\partial}{\partial t_i} $, this is a derivation of $A$.
    \[
        \mathrm{d}(f)=\sum_{i=1}^n f_i(t_1,\ldots ,t_n) \frac{\partial f}{\partial t_i}
    \]
    If $\mathrm{d}_1,\mathrm{d}_2$ are derivations of $A$, then $[\mathrm{d}_1,\mathrm{d}_2]$ is again a derivation of $A$. We can check it by the following calculation:
    \begin{align*}
        [\mathrm{d}_1,\mathrm{d}_2](ab) & =\mathrm{d}_1(\mathrm{d}_2(ab))-\mathrm{d}_2(\mathrm{d}_1(ab))                                                               \\
                                        & =\mathrm{d}_1(\mathrm{d}_2(a)b+a\mathrm{d}_2(b))-\mathrm{d}_2(\mathrm{d}_1(a)b+a\mathrm{d}_1(b))                             \\
                                        & =\mathrm{d}_1(\mathrm{d}_2(a))b+\mathrm{d}_1(a\mathrm{d}_2(b))+\mathrm{d}_2(\mathrm{d}_1(a))b+\mathrm{d}_2(a\mathrm{d}_1(b)) \\
                                        & =\mathrm{d}_1(\mathrm{d}_2(a))b+\mathrm{d}_2(\mathrm{d}_1(a))b+\mathrm{d}_1(a\mathrm{d}_2(b))+\mathrm{d}_2(a\mathrm{d}_1(b)) \\
                                        & =\mathrm{d}_1(\mathrm{d}_2(a))b+\mathrm{d}_2(\mathrm{d}_1(a))b+a[\mathrm{d}_1,\mathrm{d}_2](b)                               \\
                                        & =[\mathrm{d}_1,\mathrm{d}_2](a)b+a[\mathrm{d}_1,\mathrm{d}_2](b)                                                             \\
                                        & =[\mathrm{d}_1,\mathrm{d}_2](ab)
    \end{align*}
    Therefore $[\mathrm{d}_1,\mathrm{d}_2]$ is a derivation of $A$.
\end{example}

The above example means that:
\[
    \operatorname{Der}(A)=\{\ \mathrm{d}: A \to A \mid \mathrm{d} \text{ is a derivation} \} < \operatorname{End}_F(A)^{(-)}
\]
In particular, when $A=F[t_1,\ldots ,t_n]$, we have:
\[
    \operatorname{Der}(A)=\{\ \sum_{i=1}^n f_i(t_1,\ldots ,t_n) \frac{\partial}{\partial t_i} \mid f_i(t_1,\ldots ,t_n) \in A \} < \operatorname{End}_F(A)^{(-)}
\]
We denote this algebra by $W(n)$. It is called the $n$-th Witt algebra. It is a simple, infinite dimensional Lie algebra (when $\operatorname{char}F=0$).

Thus every element of $W(n)$ has the form
\[
    D=\sum_{i=1}^n f_i\,\partial_i,
    \qquad f_i\in A.
\]

\begin{definition}
    The divergence map is defined by
    \[
        \operatorname{div}:W(n)\to A,
        \qquad
        \operatorname{div}\!\left(\sum_{i=1}^n f_i\partial_i\right)
        =\sum_{i=1}^n \partial_i(f_i).
    \]
    Its kernel is the \emph{special} Lie algebra
    \[
        S(n):=\ker(\operatorname{div})
        =\left\{
        \sum_{i=1}^n f_i\partial_i\in W(n)
        \;\middle|\;
        \sum_{i=1}^n \partial_i(f_i)=0
        \right\}.
    \]
    Equivalently, $S(n)$ is the Lie algebra of polynomial vector fields preserving
    the standard volume form
    \[
        \omega=dx_1\wedge\cdots\wedge dx_n.
    \]
\end{definition}


A basic family of elements of $S(n)$ is given by
\[
    D_{ij}(f):=\partial_j(f)\partial_i-\partial_i(f)\partial_j,
    \qquad
    1\le i<j\le n,\quad f\in A,
\]
since
\[
    \operatorname{div}(D_{ij}(f))
    =
    \partial_i\partial_j(f)-\partial_j\partial_i(f)=0.
\]
One can check
\[
    S(n)=< D_{ij}(f) \mid 1\le i<j\le n, f\in A >_{Lie}
\]
These elements generate $S(n)$ as a Lie algebra.

\begin{definition}
    The grading of $W(n)$ is defined by
    Give $A$ its standard grading by declaring $\deg x_i=1$. Then $W(n)$ becomes a
    $\mathbb Z$-graded Lie algebra with
    \[
        \deg(x^a\partial_i)=|a|-1,
        \qquad
        |a|:=a_1+\cdots+a_n,
    \]
    so that
    \[
        W(n)=\bigoplus_{k\ge -1} W(n)_k.
    \]

    and
    \[
        W(n)_0
        =
        \operatorname{span}_F\{x_j\partial_i\mid 1\le i,j\le n\}
        \cong \mathfrak{gl}_n(F).
    \]
\end{definition}


Since the divergence map is homogeneous of degree $0$, the subalgebra $S(n)$ is
graded:
\[
    S(n)=\bigoplus_{k\ge -1} S(n)_k,
    \qquad
    S(n)_k=S(n)\cap W(n)_k.
\]
Moreover,
\[
    S(n)_{-1}=W(n)_{-1},
\]
and
\[
    S(n)_0
    =
    \left\{
    \sum_{i,j} a_{ij}x_j\partial_i
    \;\middle|\;
    \sum_i a_{ii}=0
    \right\}
    \cong \mathfrak{sl}_n(F),
\]
because
\[
    \operatorname{div}\!\left(\sum_{i,j} a_{ij}x_j\partial_i\right)
    =
    \sum_i a_{ii}.
\]

In characteristic $0$, the special algebra is perfect:
\[
    [S(n),S(n)]=S(n).
\]
Hence
\[
    S(n)^{(1)}=S(n).
\]
For $n\ge 2$, the Lie algebra $S(n)$ is simple. Therefore, in characteristic
$0$ one has
\[
    W(n)\supset S(n)=S(n)^{(1)},
\]
where both algebras are infinite-dimensional, $W(n)$ is the full Lie algebra of
polynomial vector fields, and $S(n)$ is its simple divergence-zero subalgebra.

\begin{remark}
    In characteristic $0$, the chain $W(n)\supset S(n)\supset S(n)^{(1)}$ is thus
    less informative than in positive characteristic, since the last inclusion is an
    equality.
\end{remark}

Let $A$ be an associative algebra. Fix an element $a \in A$, we define : $\mathrm{ad}(a): A \to A : b\mapsto [a,b]$. This is a derivation of $A$.

\begin{definition}
    Let $L$ be a Lie algebra. The adjoint map is defined as:
    \[
        \mathrm{ad}: L \to \operatorname{End}_F(L)^{(-)}, a \mapsto \mathrm{ad}(a)
    \]
    where $\mathrm{ad}(a): L \to L : b\mapsto [a,b]$ is an inner derivation of $L$. We also call $\mathrm{ad}(a)$ an adjoint operator.
\end{definition}

Note that $\mathrm{ad}(a) \in \operatorname{Der}(L)$ check it by the following calculation:
\begin{align*}
    \mathrm{ad}(a)([b,c]) & =\mathrm{ad}(a)(bc-cb)                       \\
                          & =[a,bc]-[a,cb]                               \\
                          & =[a,b]c+b[a,c]-[a,c]b-c[a,b]                 \\
                          & =[\mathrm{ad}(a)(b),c]+[b,\mathrm{ad}(a)(c)]
\end{align*}

"Inner" here comes from that $\mathrm{ad}(L) \triangleleft \operatorname{Der}(L)$ since $\forall \mathrm{d} \in \operatorname{Der}(L)$, we have $[\mathrm{d},\mathrm{ad}(a)]=  \mathrm{ad}(\mathrm{d}(a))$. Check it by the following calculation:
\begin{align*}
    [d,\operatorname{ad}(a)](x)
     & = d\bigl([a,x]\bigr)-[a,d(x)] \\
     & = [d(a),x]+[a,d(x)]-[a,d(x)]  \\
     & = [d(a),x]                    \\
     & = \operatorname{ad}(d(a))(x).
\end{align*}

\begin{remark}[Exponentials of nilpotent inner derivations]
    Let $L$ be a Lie algebra over a field $F$, and let
    \[
        d \in \operatorname{ad}(L) \subseteq \operatorname{Der}(L).
    \]
    Thus there exists $a \in L$ such that
    \[
        d=\operatorname{ad}(a), \qquad d(x)=[a,x] \quad \text{for all } x\in L.
    \]

    Assume that $d$ is nilpotent, so that $d^N=0$ for some integer $N\ge 1$.
    Then one defines the exponential of $d$ by
    \[
        \exp(d):=\sum_{n=0}^{\infty}\frac{d^n}{n!}.
    \]
    Since $d$ is nilpotent, this is actually a finite sum:
    \[
        \exp(d)=I+d+\frac{d^2}{2!}+\cdots+\frac{d^{N-1}}{(N-1)!}.
    \]
    Hence $\exp(d)$ is a well-defined linear endomorphism of $L$.

    Moreover, $\exp(d)$ is invertible, with inverse
    \[
        \exp(d)^{-1}=\exp(-d).
    \]
    Since $d$ is a derivation, $\exp(d)$ is in fact a Lie algebra automorphism:
    \[
        \exp(d)([x,y])=[\exp(d)(x),\exp(d)(y)]
        \qquad \text{for all } x,y\in L.
    \]
    Therefore
    \[
        \exp(d)\in \operatorname{Aut}(L).
    \]

    In particular, if $d=\operatorname{ad}(a)$, then
    \[
        \exp(\operatorname{ad}(a))(x)
        =
        x+[a,x]+\frac{1}{2!}[a,[a,x]]+\frac{1}{3!}[a,[a,[a,x]]]+\cdots.
    \]
    When $\operatorname{ad}(a)$ is nilpotent, this sum is finite. The automorphism
    \[
        \exp(\operatorname{ad}(a))
    \]
    is called an inner automorphism arising from the nilpotent inner derivation $\operatorname{ad}(a)$.

    A fundamental example is the matrix Lie algebra $L=\mathfrak{gl}_n(F)$. For $A\in \mathfrak{gl}_n(F)$, one has
    \[
        \operatorname{ad}(A)(X)=[A,X]=AX-XA,
    \]
    and
    \[
        \exp(\operatorname{ad}(A))(X)=e^A X e^{-A}.
    \]
    Thus, in the matrix case, the exponential of an inner derivation is exactly conjugation by the matrix exponential.
\end{remark}

\begin{proposition}
    Let $L$ be a Lie algebra over a field $F$, and let $d\in \operatorname{Der}(L)$ be nilpotent. Then
    \[
        \exp(d):=\sum_{n=0}^{\infty}\frac{d^n}{n!}
    \]
    is a Lie algebra automorphism of $L$.
\end{proposition}

\begin{proof}
    Since $d$ is nilpotent, there exists $N\ge 1$ such that $d^N=0$. Hence
    \[
        \exp(d)=I+d+\frac{d^2}{2!}+\cdots+\frac{d^{N-1}}{(N-1)!}
    \]
    is a finite sum, so it is a well-defined linear endomorphism of $L$.

    Moreover,
    \[
        \exp(d)\exp(-d)=\exp(-d)\exp(d)=I,
    \]
    because $d$ commutes with itself and the usual exponential identity holds for polynomials in $d$.
    Thus $\exp(d)$ is invertible, with inverse $\exp(-d)$.

    It remains to show that $\exp(d)$ preserves the Lie bracket. Define
    \[
        \phi_t:=\exp(td)=\sum_{n=0}^{\infty}\frac{t^n d^n}{n!},
    \]
    where the sum is again finite since $d$ is nilpotent. For $x,y\in L$, set
    \[
        F(t):=\phi_t([x,y]), \qquad G(t):=[\phi_t(x),\phi_t(y)].
    \]
    We show that $F(t)=G(t)$.

    First,
    \[
        F'(t)=d(\phi_t([x,y]))=\phi_t(d([x,y])).
    \]
    Since $d$ is a derivation,
    \[
        d([x,y])=[d(x),y]+[x,d(y)].
    \]
    Therefore
    \[
        F'(t)=\phi_t([d(x),y]+[x,d(y)])
        =[\phi_t(d(x)),\phi_t(y)]+[\phi_t(x),\phi_t(d(y))].
    \]
    On the other hand,
    \[
        G'(t)=[\phi_t'(x),\phi_t(y)]+[\phi_t(x),\phi_t'(y)]
        =[\phi_t(d(x)),\phi_t(y)]+[\phi_t(x),\phi_t(d(y))].
    \]
    Hence
    \[
        F'(t)=G'(t).
    \]
    Also,
    \[
        F(0)=[x,y]=G(0).
    \]
    It follows that $F(t)=G(t)$ for all $t$, and in particular at $t=1$ we get
    \[
        \exp(d)([x,y])=[\exp(d)(x),\exp(d)(y)].
    \]
    Thus $\exp(d)$ is a Lie algebra automorphism of $L$.
\end{proof}

\subsection{They come from brackets}

\begin{definition}[Poisson bracket]
    Let $A$ be a commutative associative algebra over a field $\Bbbk$.
    A \emph{Poisson bracket} on $A$ is a bilinear map
    \[
        \{\cdot,\cdot\}:A\times A\to A
    \]
    such that for all $f,g,h\in A$ the following conditions hold:
    \begin{enumerate}
        \item \textbf{Skew-symmetry:}
              \[
                  \{f,g\}=-\{g,f\}.
              \]
        \item \textbf{Jacobi identity:}
              \[
                  \{f,\{g,h\}\}+\{g,\{h,f\}\}+\{h,\{f,g\}\}=0.
              \]
        \item \textbf{Leibniz rule:}
              \[
                  \{f,gh\}=\{f,g\}h+g\{f,h\}.
              \]
    \end{enumerate}
    In other words, $(A,\{\cdot,\cdot\})$ is a Lie algebra, and for each fixed
    $f\in A$, the map
    \[
        \mathrm{ad}_f:A\to A,\qquad \mathrm{ad}_f(g)=\{f,g\},
    \]
    is a derivation of the algebra $A$.
    An algebra $A$ equipped with such a bracket is called a \emph{Poisson algebra}.
\end{definition}

\begin{remark}
    The Leibniz rule in the second variable follows automatically from skew-symmetry:
    \[
        \{fg,h\}=f\{g,h\}+g\{f,h\}.
    \]
    Thus a Poisson bracket is precisely a Lie bracket compatible with the
    commutative multiplication on $A$ by the derivation property.
\end{remark}

\begin{example}[Canonical Poisson bracket on $\mathbb{R}^{2n}$]
    Let
    \[
        A=C^\infty(\mathbb{R}^{2n}),
        \qquad
        (q_1,\dots,q_n,p_1,\dots,p_n)
    \]
    be the standard coordinates. Define
    \[
        \{f,g\}
        =
        \sum_{i=1}^n
        \left(
        \frac{\partial f}{\partial q_i}\frac{\partial g}{\partial p_i}
        -
        \frac{\partial f}{\partial p_i}\frac{\partial g}{\partial q_i}
        \right).
    \]
    Then $\{\cdot,\cdot\}$ is a Poisson bracket on $A$.
    In particular,
    \[
        \{q_i,q_j\}=0,\qquad
        \{p_i,p_j\}=0,\qquad
        \{q_i,p_j\}=\delta_{ij}.
    \]
\end{example}

\begin{definition}
    Let $F$ be a field of characteristic $0$, and let $\mathcal{A}$ be an algebra of differentiable functions in one variable $x$.
    For $f,g\in \mathcal{A}$, define
    \[
        [f,g]:=fg'-f'g,
    \]
    where $f'=\dfrac{df}{dx}$ and $g'=\dfrac{dg}{dx}$.
\end{definition}

\begin{definition}
    To each $f\in \mathcal{A}$, associate the vector field
    \[
        X_f:=f(x)\frac{d}{dx}.
    \]
\end{definition}

\begin{proposition}
    For $f,g\in \mathcal{A}$, one has
    \[
        [X_f,X_g]=X_{[f,g]}.
    \]
    Equivalently,
    \[
        [X_f,X_g]=\left(fg'-gf'\right)\frac{d}{dx}.
    \]
\end{proposition}

\begin{proof}
    For any test function $u$, we have
    \[
        X_f(u)=fu', \qquad X_g(u)=gu'.
    \]
    Thus
    \[
        X_f(X_g(u))=f(gu')'=fg'u'+fgu'',
    \]
    and similarly
    \[
        X_g(X_f(u))=g(fu')'=gf'u'+gfu''.
    \]
    Subtracting, we obtain
    \[
        [X_f,X_g](u)=X_f(X_g(u))-X_g(X_f(u))=(fg'-gf')u'.
    \]
    Therefore
    \[
        [X_f,X_g]=\left(fg'-gf'\right)\frac{d}{dx}=X_{[f,g]}.
    \]
\end{proof}

\begin{proposition}
    The operation
    \[
        [f,g]=fg'-f'g
    \]
    is bilinear and skew-symmetric.
\end{proposition}

\begin{proof}
    Bilinearity follows immediately from the linearity of differentiation. Also,
    \[
        [g,f]=gf'-g'f=-(fg'-f'g)=-[f,g].
    \]
    Hence the bracket is skew-symmetric.
\end{proof}

\begin{theorem}
    The bracket
    \[
        [f,g]=fg'-f'g
    \]
    defines a Lie algebra structure on $\mathcal{A}$.
\end{theorem}

\begin{proof}
    By the previous proposition, it remains only to verify the Jacobi identity. This follows from the fact that the map
    \[
        f\longmapsto X_f=f\frac{d}{dx}
    \]
    identifies this bracket with the usual Lie bracket of vector fields on the line, and the latter always satisfies the Jacobi identity. Hence
    \[
        [f,[g,h]]+[g,[h,f]]+[h,[f,g]]=0
    \]
    for all $f,g,h\in \mathcal{A}$.
\end{proof}

\subsection{They come from Lie groups}

\begin{definition}
    Let \(G\) be a Lie group with identity element \(e\). The \emph{Lie algebra} of \(G\), denoted by \(\operatorname{Lie}(G)\) or \(\mathfrak g\), is the tangent space
    \[
        \mathfrak g:=T_eG
    \]
    equipped with the bracket induced by left-invariant vector fields.
\end{definition}

\begin{definition}
    For each \(X\in T_eG\), the corresponding \emph{left-invariant vector field} \(\widetilde X\) on \(G\) is defined by
    \[
        \widetilde X_g:=(dL_g)_e(X),
        \qquad g\in G,
    \]
    where \(L_g:G\to G\) is left translation, \(L_g(h)=gh\).
\end{definition}

\begin{proposition}
    The assignment \(X\mapsto \widetilde X\) gives a bijection between \(T_eG\) and the space of left-invariant vector fields on \(G\).
\end{proposition}

\begin{definition}
    For \(X,Y\in \mathfrak g=T_eG\), define
    \[
        [X,Y]:=[\widetilde X,\widetilde Y](e),
    \]
    where on the right-hand side \([\widetilde X,\widetilde Y]\) is the usual Lie bracket of vector fields on \(G\).
\end{definition}

\begin{proposition}
    The bracket defined above makes \(T_eG\) into a Lie algebra. In particular, it is bilinear, alternating, and satisfies the Jacobi identity.
\end{proposition}

\begin{example}
    If \(G=GL_n(F)\), where \(F=\mathbb R\) or \(\mathbb C\), then
    \[
        \operatorname{Lie}(G)=T_IGL_n(F)\cong M_n(F),
    \]
    and the Lie bracket is
    \[
        [X,Y]=XY-YX.
    \]
    Hence
    \[
        \operatorname{Lie}(GL_n(F))=\mathfrak{gl}_n(F).
    \]
\end{example}

\begin{example}
    If \(G=SL_n(F)\), then
    \[
        \operatorname{Lie}(SL_n(F))=\mathfrak{sl}_n(F)
        =\{X\in M_n(F)\mid \operatorname{tr}(X)=0\}.
    \]
    Indeed, differentiating the defining equation \(\det(g)=1\) at the identity gives the trace-zero condition.
\end{example}

\begin{example}
    If \(G=O_n(\mathbb R)\), then
    \[
        \operatorname{Lie}(O_n(\mathbb R))=\mathfrak{so}_n(\mathbb R)
        =\{X\in M_n(\mathbb R)\mid X^T=-X\}.
    \]
    This follows by differentiating the relation \(g^Tg=I\) at the identity.
\end{example}

